\section{Validation, reproducibility, and further research}
\label{sec:validation}

The composition suite checks 596 complete-process expectations and ten
targeted mutation observations. It covers direct computations, repaired
zero computations, nested binders, contextual induction, transport,
scope failures, false conclusions, scalar-major aliases and malformed
JSON, under ordinary Python and \texttt{-O}. The mutation controls exercise
major verification, equation-only induction, closed-goal scope, genuine
successor construction and lexical-floor transmission. Finite testing is
evidence about these implementations, not a proof of the universal
metatheorems.

One historical raw alias has a documented interface difference: the
545-byte guarded comparator rejects a null-major packet when its affine
dispatch attempts \texttt{None*3}; its translated empty-list major is a
valid accepted zero-instance packet. Canonical completeness does not
promise preservation of every redundant raw spelling.

The independent arithmetic diagnostics evaluate the finite-table and
capped-clock constructions on standard finite inputs. They are kept
separate from kernel certificates and never described as proving the
claims in nonstandard models. The re-examined progression-binder prototype
is outside the selected checker comparison: its theory is already present
in the source corpus, and its 576-byte implementation is larger. It is not
a newly proposed replacement for the measured calculus.

Run \texttt{python3 reproduce.py} from a fresh extraction to remeasure
the selected modules and frozen corpus, run the public checker suites,
and repeat the arithmetic diagnostics. Run \texttt{python3 build.py}
to rebuild the paper with PDFLaTeX. No network access, external proof
oracle, or shell escape is required. The standalone expanded \TeX\ source
can also be compiled on its own. Rebuilding can change PDF metadata, so
the delivered hash manifest should be checked before rebuilding.

\paragraph{Further research questions.}
The most concrete remaining mathematical question is whether the
self-seeded precursor can upgrade its selected right predecessor to a
left predecessor. A positive proof of $P(Sn)=n$, or any other internal
left inverse, would remove the injectivity assumption from the known
completion. A negative answer needs a model of the complete precursor,
not merely an acyclic noninjective successor graph.

The finite-cycle proof uses a different fixed term for each external
cycle length. A uniform internal finite-cycle construction could provide
stronger cancellation information, but cannot be obtained by silently
adding an evaluator for all descriptions. The existing all-starts theorem
already combines the progression binder's Step and Base interfaces. The
remaining question there is whether that fusion admits a smaller complete
implementation; re-proving the fusion alone would repeat the prior work.
Finally, a compact Lean or Rocq formalization of the sorted opener,
lexical scope invariant, arithmetic bootstrap and both interpretation
composites would provide much stronger assurance than the present
written proofs and finite tests.
